\section{Introduction}
\label{sec:introduction}

Financial transaction systems process billions of electronic operations daily under constant exposure to adversarial fraud syndicates \cite{leborgne2022handbook}. In retail banking, automated card fraud triage faces three severe operational constraints that violate standard supervised learning assumptions:
\begin{enumerate}
    \item \textbf{Extreme Class Imbalance:} Fraudulent events represent an ultra-rare class ($<0.2\%$). In the benchmark European credit card stream \cite{dalpozzolo2018credit}, only $492$ out of $284,807$ transactions ($0.1727\%$) are fraudulent. Trivial majority-class prediction yields $99.83\%$ accuracy while catching zero financial losses.
    \item \textbf{Verification Lag:} Ground-truth fraud labels are not available at transaction authorization time, emerging only days or weeks later following cardholder dispute resolutions \cite{carcillo2018scarff}.
    \item \textbf{Adversarial Non-Stationarity:} Fraud strategies evolve dynamically, while legitimate consumer spending undergoes circadian and seasonal shifts, creating ongoing distribution drift \cite{dalpozzolo2015calibrating}.
\end{enumerate}

To address label scarcity and delayed feedback, fraud triage can be formulated as an \emph{unsupervised anomaly detection} problem: models learn the probability density or bounding envelope of legitimate cardholder behavior, scoring unseen transactions by their deviation from normal patterns.

\subsection{Phase 1 Scope and Contributions}
In Phase~1 of Project T033, we evaluate classical machine learning and statistical density estimation methods for credit card fraud triage. We analyze Gaussian Mixture Models (GMM) estimated via Expectation-Maximization (EM) \cite{dempster1977maximum} alongside four classical baselines: PCA Reconstruction Error \cite{shyu2003novel}, Local Outlier Factor (LOF) \cite{breunig2000lof}, One-Class SVM (OC-SVM) \cite{scholkopf2001estimating}, and Isolation Forest \cite{liu2008isolation}. Deep autoencoder architectures are reserved for Phase~2.

The principal contributions of this work are:
\begin{itemize}
    \item \textbf{Strict Chronological Protocol:} A non-overlapping temporal split (60\% train, 20\% validation, 20\% test) spanning 48 hours with all scalers fit strictly on legitimate training data to eliminate data leakage.
    \item \textbf{Gaussianity and Distributional Auditing:} Empirical statistical testing of normality across PCA features, identifying severe leptokurtosis that challenges parametric density assumptions.
    \item \textbf{Capacity-Constrained Evaluation:} Evaluating models under realistic operational constraints using PR-AUC, Recall at fixed False Positive Rates ($\text{FPR} \le 0.5\%$), and Precision at investigator review capacity.
    \item \textbf{Systematic Model Selection:} Comparing GMM covariance topologies and component orders ($K$) governed by AIC/BIC, highlighting the divergence between density likelihood and fraud triage contrast.
\end{itemize}
